\documentclass[sigconf, nonacm]{acmart}

% --- PACKAGES ---
\usepackage{subcaption}
\usepackage{multirow}
\usepackage{colortbl}
\usepackage{enumitem}

% --- MAKE LINKS CLICKABLE AND BLUE ---
\hypersetup{
    colorlinks=true,
    linkcolor=blue,
    filecolor=magenta,
    urlcolor=blue,
    citecolor=blue,
}

% Metadata
\copyrightyear{2025}
\acmYear{2025}
\setcopyright{none}
\acmConference[]{}{}

\settopmatter{printacmref=false, printfolios=true}
\pagestyle{plain}

\begin{document}

% --- TITLE & AUTHOR ---
\title{Benchmarking Geospatial Foundation Models for
Agriculture Applications}

\author{Sanmay Das}
\affiliation{%
  \institution{University of California, Riverside}
  \city{Riverside}
  \country{USA}
}
\email{sdas050@ucr.edu}

\maketitle

% ============================================================
\section{Introduction}
% ============================================================

Accurate crop type mapping at scale underpins food security
monitoring, yield estimation, and agricultural
decision-making across the U.S., where 880 million acres of
farmland---nearly 39\% of all land---cannot be monitored
through manual surveys alone. Geospatial Foundation Models
(GeoFMs) such as Prithvi~\cite{jakubik2023foundation},
SatMAE~\cite{cong2022satmae}, and
SpectralGPT~\cite{hong2024spectralgpt}---pretrained on
massive Earth observation corpora---promise generalizable
representations for multi-temporal crop segmentation and
change detection with minimal task-specific supervision.
Yet a fundamental question remains: \textit{do these models
generalize across geographies, or do reported results merely
reflect dataset idiosyncrasies?}

Existing benchmarks leave this question unresolved.
PASTIS~\cite{garnot2021panoptic},
BreizhCrop~\cite{russwurm2019breizhcrops},
Sen4AgriNet~\cite{sykas2021sen4agrinet}, and
EuroCropsML~\cite{reuss2025eurocropsml} are
all European in scope, with labels sourced from
farmer self-declarations that introduce systematic noise.
GEO-Bench~\cite{lacoste2023geobench} randomly subsamples
to 20{,}000 chips, preventing assessment of spatial
autocorrelation and geographic generalization.
PANGAEA~\cite{marsocci2024pangaea} finds that supervised
baselines outperform GeoFMs under distribution shift,
while PhilEO Bench~\cite{fibaek2024phileo} identifies a
detail gap in which foundation models lose fine-grained
spatial features on dense prediction tasks---yet neither
targets U.S.\ agricultural conditions or provides
controlled multi-state comparisons.

We present a systematic benchmark evaluating Prithvi,
SpectralGPT, and SatMAE on multi-temporal crop segmentation
and change detection using raw Sentinel-2
imagery~\cite{sentinel2_eoportal} across four climatically
diverse U.S.\ states: Iowa, North Carolina, California, and
Minnesota. Under controlled, consistent conditions---identical
preprocessing, contiguous 800--1{,}300~km$^{2}$ study
regions, and strict geographic train/val/test splits---our
results demonstrate that \textbf{spectral encoding provides
a more generalizable inductive bias than temporal encoding}
for cross-regional crop mapping, and expose systematic
failure modes including majority class collapse and
minority crop invisibility under distribution shift.

% ============================================================
\section{Dataset \& Methodology}
% ============================================================

\textbf{Study Regions.}
We selected three spatially contiguous 800--1{,}300~km$^{2}$
regions per state assigned to train, validation, and test
splits with no spatial overlap, covering Iowa (Ames /
Cedar~Rapids / Sheldon--Sibley), North Carolina
(Rocky~Mount / Clinton / Belhaven), California
(Rio~Vista / Merced / Bakersfield), and Minnesota
(Willmar / Crookston / Rochester). States were chosen to
maximize climatic and agronomic diversity: Iowa and Minnesota
represent the annual row-crop Corn Belt; North Carolina
introduces coastal ecotones and organic muck soils; and
California introduces perennial orchards and arid irrigated
agriculture under the most severe distribution shift.

\textbf{Imagery \& Labels.}
Sentinel-2 Level-2A imagery was acquired via the Copernicus
Data Space Ecosystem~\cite{cdse_annual_report} for the 2024
growing season (March--September) using six bands: B2, B3,
B4, B8A, B11, B12~\cite{sentinel2_bands}. Three phenological
composites (early / mid / late season) were formed by
selecting the least-cloudy scene per window using the SCL
mask. All scenes were reprojected to EPSG:5070 at 10~m and
stacked into 18-band GeoTIFFs. Chips were extracted with
50\% sliding-window overlap and a 5\% per-chip cloud
threshold at model-native resolutions: 224$\times$224
(Prithvi), 128$\times$128 (SpectralGPT), and 96$\times$96
(SatMAE). Ground truth was derived from the USDA Cropland
Data Layer (CDL)~\cite{usda_cdl} for 2024, mapped to 13
classes: Natural Vegetation, Forest, Corn, Soybeans,
Wetlands, Developed/Barren, Open Water, Winter Wheat,
Alfalfa, Fallow/Idle, Cotton, Sorghum, and Other.

\textbf{Models.}
Candidate selection required native multispectral support
beyond RGB---eliminating Scale-MAE~\cite{reed2023scale}
(RGB-only pretraining) and SAM~\cite{osco2023segment}
(natural image pretraining). Three models were retained:

\begin{itemize}[leftmargin=*, itemsep=1pt, topsep=2pt]
  \item \textbf{Prithvi}~\cite{jakubik2023foundation}: ViT
    masked autoencoder pretrained on multi-temporal HLS
    imagery. Uses 3D spatiotemporal patch embeddings and
    lacks explicit geolocation embeddings. Segmentation head:
    4-layer ConvTranspose2D neck + FCN. Change detection:
    Siamese encoder + FCN. Trained 80 epochs, Adam,
    LR$=1.5\!\times\!10^{-5}$, weighted cross-entropy.
  \item \textbf{SpectralGPT}~\cite{hong2024spectralgpt}:
    Generative model with 3D spatial-spectral coupling,
    adapted for multi-temporal input. Segmentation head: MLP
    + UPerNet. Change detection: Siamese encoder + FPN.
    Trained 250 epochs, AdamW, LR$=1\!\times\!10^{-4}$,
    cross-entropy.
  \item \textbf{SatMAE}~\cite{cong2022satmae}: GroupChannels
    ViT-Large masked autoencoder with per-group spectral
    patch embeddings, pretrained on fMoW-Sentinel.
    Segmentation head: FPN. Change detection: Siamese encoder
    + FCN. Trained 100 epochs, AdamW, cross-entropy.
\end{itemize}

\textbf{Inference \& Evaluation.}
Full-scene segmentation inference followed the TerraTorch
tiled protocol~\cite{gomes2025terratorch}: reflect-padded
sliding window, cosine-tapered blend weights, and an
8-pixel border discard to suppress edge artifacts.
For change detection, single mid-season scenes from 2023
(T1) and 2024 (T2) serve as the bi-temporal image pair.
Change masks are derived by directly comparing 2023 and
2024 CDL annotations, both remapped to the same 13-class
schema: any pixel where the land cover class differs
between years is labeled changed (1), same-class pixels
as unchanged (0), and pixels with missing CDL data in
either year are excluded (255). Both T1 and T2 are
cloud-filtered using their respective SCL masks, requiring
$\geq$95\% valid pixels per chip. All three models use a
Siamese encoder with shared weights; Prithvi and SatMAE
use an FCN head, SpectralGPT uses an FPN head.
Segmentation was evaluated with mIoU and per-class
IoU; change detection with Precision, Recall, F1, and
IoU\textsubscript{change} following the SpectralGPT
protocol~\cite{hong2024spectralgpt}. All experiments ran
on a single NVIDIA A100 GPU.

% ============================================================
\section{Results}
% ============================================================

\subsection{Segmentation}

Table~\ref{tab:iou_results} presents per-class IoU for
Prithvi, SpectralGPT, and SatMAE across all four states;
Figure~\ref{fig:qual_all} shows the corresponding
qualitative prediction maps.

\begin{table*}[t]
\centering
\caption{Per-class IoU (\%) for Prithvi, SpectralGPT, and SatMAE across
four U.S.\ states. mIoU is averaged across all states.
NaN = class absent in that region. -- = results pending.}
\label{tab:iou_results}
\resizebox{\textwidth}{!}{%
\begin{tabular}{lcccccccccccccccc}
\toprule
\multirow{2}{*}{\textbf{Class}} &
\multicolumn{3}{c}{\textbf{Iowa}} &
\multicolumn{3}{c}{\textbf{North Carolina}} &
\multicolumn{3}{c}{\textbf{California}} &
\multicolumn{3}{c}{\textbf{Minnesota}} &
\multicolumn{3}{c}{\textbf{mIoU}} \\
\cmidrule(lr){2-4}\cmidrule(lr){5-7}
\cmidrule(lr){8-10}\cmidrule(lr){11-13}\cmidrule(lr){14-16}
& \textbf{Prithvi} & \textbf{SGPT} & \textbf{SatMAE}
& \textbf{Prithvi} & \textbf{SGPT} & \textbf{SatMAE}
& \textbf{Prithvi} & \textbf{SGPT} & \textbf{SatMAE}
& \textbf{Prithvi} & \textbf{SGPT} & \textbf{SatMAE}
& \textbf{Prithvi} & \textbf{SGPT} & \textbf{SatMAE} \\
\midrule
Natural Vegetation & 0.00  & 41.86 & -- & 11.10 & 32.06 & -- & 1.46  & 74.86 & -- & 0.01  & 14.94 & -- & 3.14  & 40.93 & -- \\
Forest             & 0.73  & 8.62  & -- & 65.42 & 88.43 & -- & 12.98 & 2.74  & -- & 2.16  & 44.38 & -- & 20.32 & 36.04 & -- \\
Corn               & 33.69 & 7.99  & -- & 9.58  & 48.83 & -- & 0.00  & 0.75  & -- & 4.04  & 46.86 & -- & 11.83 & 26.11 & -- \\
Soybean            & 38.41 & 26.70 & -- & 6.12  & 59.84 & -- & NaN   & NaN   & -- & 17.27 & 62.14 & -- & 20.60 & 49.56 & -- \\
Wetlands           & 0.00  & 2.03  & -- & 0.00  & 0.02  & -- & 0.00  & 0.45  & -- & 1.69  & 14.08 & -- & 0.42  & 4.15  & -- \\
Developed/Barren   & 1.51  & 53.35 & -- & 3.22  & 56.43 & -- & 6.95  & 23.91 & -- & 0.32  & 52.83 & -- & 3.00  & 46.63 & -- \\
Open Water         & 6.86  & 39.44 & -- & 0.13  & 90.26 & -- & 11.17 & 34.76 & -- & 0.76  & 75.35 & -- & 4.73  & 59.95 & -- \\
Winter Wheat       & 0.00  & 0.00  & -- & 0.00  & 0.87  & -- & 0.00  & 0.01  & -- & 0.00  & 0.00  & -- & 0.00  & 0.22  & -- \\
Alfalfa            & 0.00  & 30.62 & -- & 0.00  & 0.00  & -- & 0.09  & 41.08 & -- & 0.12  & 9.31  & -- & 0.05  & 20.25 & -- \\
Fallow/Idle        & 0.00  & 0.34  & -- & 0.00  & 0.25  & -- & 0.00  & 2.80  & -- & 0.00  & 23.57 & -- & 0.00  & 6.74  & -- \\
Cotton             & NaN   & NaN   & -- & 2.35  & 66.90 & -- & 0.00  & 0.00  & -- & NaN   & NaN   & -- & 1.18  & 33.45 & -- \\
Sorghum            & 0.00  & 0.00  & -- & 0.00  & 4.94  & -- & 0.00  & 0.03  & -- & 0.00  & 0.00  & -- & 0.00  & 1.24  & -- \\
Other              & 0.00  & 0.08  & -- & 4.51  & 64.75 & -- & 20.35 & 62.96 & -- & 0.93  & 8.35  & -- & 6.45  & 34.04 & -- \\
\bottomrule
\end{tabular}}
\end{table*}

SpectralGPT consistently outperforms Prithvi across all
states, achieving substantially higher mIoU on Open Water
(59.95 vs.\ 4.73), Developed/Barren (46.63 vs.\ 3.00), and
Soybean (49.56 vs.\ 20.60). Prithvi is competitive only on
Corn in Iowa (33.69 vs.\ 7.99) and Forest in North Carolina
(65.42 vs.\ 88.43)---classes where its training distribution
most closely matches the test region. Both models fail on
spectrally ambiguous minority classes: Winter Wheat
($\leq$0.87), Sorghum ($\leq$4.94), and Fallow/Idle
($\leq$0.34).

\textbf{Majority class collapse} is Prithvi's dominant
failure mode: after fine-tuning it shifts its global
prediction to the most frequent class in each training
region (Corn in Iowa, Developed/Barren in California),
absorbing all minority structures into the dominant
background. \textbf{Geographic generalization} further
separates the models. Prithvi's Iowa Corn IoU of 33.69
drops to 9.58 in North Carolina and 0.00 in California:
the latitudinal phenology shift in Minnesota causes its
temporal embeddings to fail (Corn: 4.04); the dark organic
muck soils of coastal North Carolina distort the green-up
baseline its 3D patch embeddings rely on; and Kern County's
perennial orchards---whose canopy remains spectrally
static year-round---are invisible to a sequence-dependent
temporal tracker. SpectralGPT's 3D spatial-spectral
coupling and multi-target reconstruction loss identify
materials through their radiometric fingerprint regardless
of temporal trajectory, yielding stronger cross-regional
transfer---though its Forest IoU collapses to 2.74 in
California, where forested regions have a markedly different
spectral profile from Eastern states.

\begin{figure*}[t]
\centering
\setlength{\tabcolsep}{2pt}
\newcommand{\qh}{2.6cm}% per-image height — reduce if paper overflows
\begin{tabular}{cccc}
  % ── column headers ──────────────────────────────────────────
  & \textbf{Ground Truth} & \textbf{Prithvi} & \textbf{SpectralGPT} \\[2pt]
  % ── Iowa ────────────────────────────────────────────────────
  \rotatebox{90}{\small\textbf{Iowa}} &
  \includegraphics[height=\qh,keepaspectratio]{NWIA_gt_colored.png} &
  \includegraphics[height=\qh,keepaspectratio]{NWIA_Prithvi_colored.png} &
  \includegraphics[height=\qh,keepaspectratio]{NWIA_SpectralGPT_colored.png} \\[2pt]
  % ── North Carolina ───────────────────────────────────────────
  \rotatebox{90}{\small\textbf{N.\ Carolina}} &
  \includegraphics[height=\qh,keepaspectratio]{EastNC_gt_colored.png} &
  \includegraphics[height=\qh,keepaspectratio]{EastNC_Prithvi_colored.png} &
  \includegraphics[height=\qh,keepaspectratio]{EastNC_SpectralGPT_colored.png} \\[2pt]
  % ── California ───────────────────────────────────────────────
  \rotatebox{90}{\small\textbf{California}} &
  \includegraphics[height=\qh,keepaspectratio]{SouthCA_gt_colored.png} &
  \includegraphics[height=\qh,keepaspectratio]{SouthCA_Prithvi_colored.png} &
  \includegraphics[height=\qh,keepaspectratio]{SouthCA_SpectralGPT_colored.png} \\[2pt]
  % ── Minnesota ────────────────────────────────────────────────
  \rotatebox{90}{\small\textbf{Minnesota}} &
  \includegraphics[height=\qh,keepaspectratio]{SouthMN_gt_colored.png} &
  \includegraphics[height=\qh,keepaspectratio]{SouthMN_Prithvi_colored.png} &
  \includegraphics[height=\qh,keepaspectratio]{SouthMN_SpectralGPT_colored.png} \\
\end{tabular}
\caption{Qualitative segmentation maps across all four test regions.
  \textit{Iowa}: Prithvi collapses to Corn (yellow); SpectralGPT recovers
  Developed/Barren, Open Water, and Natural Vegetation.
  \textit{North Carolina}: Prithvi collapses to Natural Vegetation on
  coastal muck soils; SpectralGPT recovers Open Water (90.26\%) and Forest
  (88.43\%).
  \textit{California}: Prithvi collapses to Developed/Barren in Kern County
  orchards; SpectralGPT identifies Natural Vegetation (74.86\%) and Alfalfa
  (41.08\%).
  \textit{Minnesota}: Prithvi fails on Corn (4.04\%) due to latitudinal
  phenology shift; SpectralGPT achieves Corn (46.86\%) and Soybean
  (62.14\%).}
\label{fig:qual_all}
\end{figure*}

\subsection{Change Detection}

\begin{table}[h]
\centering
\caption{Change detection on Iowa (NWIA). F1 and
IoU\textsubscript{ch} are primary metrics.
Train: CentIA, Val: EastIA, Test: NWIA.}
\label{tab:cd_results}
\begin{tabular}{lcccc}
\toprule
\textbf{Model} & \textbf{Prec.} & \textbf{Rec.}
& \textbf{F1} & \textbf{IoU\textsubscript{ch}} \\
\midrule
SpectralGPT & 92.09 & 54.16 & 68.20 & 51.75 \\
Prithvi     & 82.53 & 41.28 & 55.03 & 37.96 \\
SatMAE      & 70.42 & 29.01 & 41.10 & 25.86 \\
\bottomrule
\end{tabular}
\end{table}

The model ranking (SpectralGPT $>$ Prithvi $>$ SatMAE)
mirrors segmentation, confirming that spectral encoding
generalizes across both dense prediction tasks. All models
show high precision but low recall (29--54\%), consistent
with the class imbalance in the change mask where genuine
land cover transitions are rare under the coarse CDL
remapping. SatMAE's lower recall (29.01) reflects the
GroupChannels ViT's group-averaging, which dilutes the
fine-grained bi-temporal difference signal. Prithvi's
explicit temporal embeddings provide a modest advantage
over SatMAE in bi-temporal change detection despite
failing under geographic shift in segmentation.

\begin{figure}[h]
\centering
\begin{subfigure}{0.22\columnwidth}
    \includegraphics[width=\linewidth,height=3.5cm,
    keepaspectratio]{NWIA_CD_GT_colored.png}
    \caption*{GT}
\end{subfigure}
\hfill
\begin{subfigure}{0.22\columnwidth}
    \includegraphics[width=\linewidth,height=3.5cm,
    keepaspectratio]{NWIA_CD_Prithvi_colored.png}
    \caption*{Prithvi}
\end{subfigure}
\hfill
\begin{subfigure}{0.22\columnwidth}
    \includegraphics[width=\linewidth,height=3.5cm,
    keepaspectratio]{NWIA_CD_SatMAE_colored.png}
    \caption*{SatMAE}
\end{subfigure}
\hfill
\begin{subfigure}{0.22\columnwidth}
    \includegraphics[width=\linewidth,height=3.5cm,
    keepaspectratio]{NWIA_CD_SpectralGPT_colored.png}
    \caption*{SGPT}
\end{subfigure}
\caption{Iowa (NWIA) change detection maps. White =
detected change; black = no change. SpectralGPT most
closely recovers the spatial structure of the GT;
SatMAE produces the sparsest, noisiest change map.}
\label{fig:qual_cd}
\end{figure}

% ============================================================
\section{Conclusion}
% ============================================================

We present the first controlled, multi-state benchmark
comparing Prithvi, SpectralGPT, and SatMAE for multi-temporal
crop segmentation and change detection across climatically
diverse U.S.\ agricultural regions. Spectral encoding
(SpectralGPT) provides substantially stronger geographic
generalization than temporal encoding (Prithvi): Prithvi
collapses to majority class prediction under phenological
mismatch, soil reflectance shifts, and perennial crop
environments, while SpectralGPT's radiometric fingerprinting
transfers across geographies with no temporal assumptions.
Both models fail consistently on spectrally ambiguous minority
classes (Winter Wheat, Sorghum, Fallow/Idle), and all three
models under-detect change due to inherent class imbalance
in the binary change mask. These findings reveal that no
current GeoFM achieves robust performance across the full
13-class schema under geographic distribution shift,
motivating architectural designs that combine spectral
fidelity with explicit geolocation awareness.
Future work will extend the benchmark to 8 additional
U.S.\ states and evaluate Prithvi-EO-2.0's geolocation
embeddings as a remedy for the region-locking behavior
identified here.

\bibliographystyle{ACM-Reference-Format}
\bibliography{references}

\end{document}
